% =====================================================================
%  Module 2 : Sécurité et gestion des accès
% =====================================================================
\section{Sécurité et gestion des accès}

\begin{frame}{IAM : Identity and Access Management}
  \begin{itemize}
    \setlength{\itemsep}{.55em}
    \item Contrôle l'accès aux services et ressources AWS
    \item \textbf{Authentification} : qui envoie la requête ?
    \item \textbf{Autorisation} : cette identité a-t-elle le droit de faire cette action sur cette ressource ?
    \item S'applique à toutes les requêtes : console, CLI, SDK
    \item Service global (pas de région), gratuit
  \end{itemize}
  \lien{https://docs.aws.amazon.com/IAM/latest/UserGuide/introduction.html}
\end{frame}

\begin{frame}{Vocabulaire}
  \renewcommand{\arraystretch}{1.45}
  \begin{tabularx}{\textwidth}{@{}>{\raggedright\arraybackslash}p{3.2cm}>{\raggedright\arraybackslash}X@{}}
    \toprule
    \textbf{Principal} & entité qui envoie une requête : utilisateur racine, utilisateur IAM, rôle, service AWS \\
    \textbf{Utilisateur IAM} & identité permanente d'une personne ou d'une application ; mot de passe et/ou clés d'accès \\
    \textbf{Groupe} & ensemble d'utilisateurs qui partagent les mêmes stratégies ; ne peut pas se connecter \\
    \textbf{Rôle} & identité sans identifiants permanents, endossée temporairement \\
    \textbf{Stratégie} (\emph{policy}) & document JSON qui autorise ou interdit des actions sur des ressources \\ \bottomrule
  \end{tabularx}
\end{frame}

\begin{frame}{Utilisateur racine}
  \begin{itemize}
    \setlength{\itemsep}{.6em}
    \item Créé avec le compte, identifié par l'adresse électronique d'inscription
    \item Tous les droits, non restreignables par IAM
    \item MFA obligatoire pour tous les types de comptes depuis juin 2025
    \item À réserver aux tâches qui l'exigent (fermeture du compte, plan de support...)
    \item Pas de clé d'accès pour l'utilisateur racine
  \end{itemize}
  \bigskip
  Dans ce cours : accès par un utilisateur IAM fourni, jamais par le compte racine
  \source{AWS, « AWS IAM now enforces MFA for root users across all account types », 2025.}
\end{frame}

\begin{frame}{Utilisateurs, groupes, rôles}
  \centering
  \includegraphics[width=.95\textwidth]{iam-identites}
\end{frame}

\begin{frame}{Rôles}
  \begin{itemize}
    \setlength{\itemsep}{.5em}
    \item Pas de mot de passe ni de clé permanente
    \item Endossé par une entité autorisée : service AWS (EC2, Lambda), utilisateur, autre compte, utilisateur fédéré
    \item Identifiants temporaires délivrés par AWS STS, renouvelés automatiquement
    \item Deux stratégies :
    \begin{itemize}
      \item \textbf{stratégie de confiance} : qui peut endosser le rôle
      \item \textbf{stratégie de permissions} : ce que le rôle autorise
    \end{itemize}
    \item Instance EC2 : rôle associé via un \emph{instance profile}
  \end{itemize}
  \bigskip
  \attention Une application sur AWS n'a jamais besoin de clé d'accès : elle utilise un rôle
\end{frame}

\begin{frame}{Structure d'une stratégie}
  \centering
  \includegraphics[width=.96\textwidth]{politique-annotee}
\end{frame}

\begin{frame}[fragile]{Stratégie avec condition}
  Exemple : interdire toute action EC2 en dehors de la région de Paris
\begin{lstlisting}[language=json]
{
  "Version": "2012-10-17",
  "Statement": [{
    "Sid": "UniquementParis",
    "Effect": "Deny",
    "Action": "ec2:*",
    "Resource": "*",
    "Condition": {
      "StringNotEquals": { "aws:RequestedRegion": "eu-west-3" }
    }
  }]
}
\end{lstlisting}
  Autres clés de condition courantes : \texttt{aws:SourceIp}, \texttt{aws:MultiFactorAuthPresent}, \texttt{aws:ResourceTag/...}
\end{frame}

\begin{frame}{Types de stratégies}
  \renewcommand{\arraystretch}{1.4}
  \begin{tabularx}{\textwidth}{@{}>{\raggedright\arraybackslash}p{4.6cm}>{\raggedright\arraybackslash}X@{}}
    \toprule
    \textbf{Gérée par AWS} & prête à l'emploi (\texttt{AdministratorAccess}, \texttt{AmazonS3ReadOnlyAccess}) ; souvent trop large \\
    \textbf{Gérée par le client} & écrite par le client, réutilisable ; adaptée au moindre privilège \\
    \textbf{En ligne} (\emph{inline}) & intégrée à une seule identité, supprimée avec elle \\
    \textbf{Basée sur la ressource} & attachée à une ressource : stratégie de bucket S3, de file SQS \\
    \textbf{SCP}, limite de permissions & plafonnent les droits possibles (organisation, utilisateur) \\ \bottomrule
  \end{tabularx}
\end{frame}

\begin{frame}{Logique d'évaluation}
  \begin{columns}[c]
    \begin{column}{.45\textwidth}
      \centering
      \includegraphics[height=.64\textheight]{evaluation-iam}
    \end{column}
    \begin{column}{.52\textwidth}
      \begin{itemize}
        \setlength{\itemsep}{.6em}
        \item Par défaut : tout est refusé
        \item Un \texttt{Allow} explicite lève le refus
        \item Un \texttt{Deny} explicite l'emporte toujours
        \item Toutes les stratégies applicables sont évaluées ensemble
      \end{itemize}
    \end{column}
  \end{columns}
  \lien{https://docs.aws.amazon.com/IAM/latest/UserGuide/reference_policies_evaluation-logic.html}
\end{frame}

\begin{frame}{ARN : Amazon Resource Name}
  \begin{itemize}
    \item Identifiant unique de toute ressource AWS
    \item Format : \texttt{arn:aws:\emph{service}:\emph{région}:\emph{compte}:\emph{ressource}}
  \end{itemize}
  \medskip
  \renewcommand{\arraystretch}{1.35}
  \begin{tabularx}{\textwidth}{@{}l>{\raggedright\arraybackslash}X@{}}
    \toprule
    \texttt{arn:aws:iam::123456789012:user/camille} & IAM : global, pas de région \\
    \texttt{arn:aws:ec2:eu-west-3:123456789012:instance/i-0abc} & une instance à Paris \\
    \texttt{arn:aws:s3:::galerie-camille} & un bucket : ni région ni compte \\
    \texttt{arn:aws:s3:::galerie-camille/*} & les objets du bucket \\ \bottomrule
  \end{tabularx}
\end{frame}

\begin{frame}{Clés d'accès et identifiants temporaires}
  \centering
  \includegraphics[width=.7\textwidth]{cle-vs-role}\par\medskip
  \raggedright
  \begin{itemize}
    \item Clé permanente (\texttt{AKIA...}) : valable jusqu'à sa désactivation
    \item Identifiant temporaire (\texttt{ASIA...}) : expire après quelques heures
    \item Clé publiée sur GitHub : mise en quarantaine par AWS en une dizaine de secondes
  \end{itemize}
  \source{Unit 42, « From Exposure to Lockdown » ; Meli et al., « How Bad Can It Git? », NDSS 2019.}
\end{frame}

\begin{frame}{Bonnes pratiques IAM}
  \begin{itemize}
    \setlength{\itemsep}{.45em}
    \item Moindre privilège : uniquement les droits nécessaires à la tâche
    \item Droits attribués aux groupes, pas aux utilisateurs
    \item Rôles et identifiants temporaires plutôt que clés d'accès
    \item MFA pour toutes les personnes
    \item Fédération d'identités (IAM Identity Center) pour les personnes en entreprise
    \item Suppression des utilisateurs, clés et droits inutilisés
    \item Comptes séparés par environnement (développement, production)
  \end{itemize}
  \bigskip
  \attention Ne jamais versionner de clé d'accès dans un dépôt Git
  \lien{https://docs.aws.amazon.com/IAM/latest/UserGuide/best-practices.html}
\end{frame}

\begin{frame}{Outils IAM}
  \renewcommand{\arraystretch}{1.45}
  \begin{tabularx}{\textwidth}{@{}>{\raggedright\arraybackslash}p{4.6cm}>{\raggedright\arraybackslash}X@{}}
    \toprule
    \textbf{IAM Policy Simulator} & tester une stratégie sans exécuter d'action \\
    \textbf{IAM Access Analyzer} & détecter les accès publics ou externes, générer des stratégies à partir de l'activité \\
    \textbf{Dernier accès} (\emph{last accessed}) & services utilisés par une identité, droits superflus \\
    \textbf{CloudTrail} & journal des appels d'API : qui, quoi, quand, d'où (90 jours gratuits) \\ \bottomrule
  \end{tabularx}
\end{frame}

% ---- réseau ---------------------------------------------------------------
\begin{frame}{VPC : Virtual Private Cloud}
  \begin{itemize}
    \setlength{\itemsep}{.5em}
    \item Réseau virtuel privé et isolé, propre au compte, dans une région
    \item Plage d'adresses en notation CIDR, par exemple \texttt{10.0.0.0/16}
    \item Contient les sous-réseaux, tables de routage, passerelles
    \item Socle de toutes les ressources réseau : EC2, RDS, ELB...
    \item Plusieurs VPC par compte (5 par région par défaut)
  \end{itemize}
  \bigskip
  \textbf{VPC par défaut} : créé dans chaque région, \texttt{172.31.0.0/16}, un sous-réseau public par AZ, passerelle Internet déjà attachée
\end{frame}

\begin{frame}{Sous-réseaux et accès Internet}
  \begin{itemize}
    \setlength{\itemsep}{.45em}
    \item Sous-réseau : plage d'adresses du VPC, dans \textbf{une seule} AZ (ex. \texttt{10.0.1.0/24})
    \item \textbf{Public} : route vers une passerelle Internet
    \item \textbf{Privé} : pas d'accès entrant depuis Internet
    \item \textbf{Internet Gateway} : relie le VPC à Internet, gratuite
    \item \textbf{NAT Gateway} : sortie vers Internet pour les sous-réseaux privés ; environ 0,045 \$/h + 0,045 \$/Go traité (\texttt{us-east-1})
    \item Adresse IPv4 publique : facturée 0,005 \$/h
  \end{itemize}
  \source{AWS, \emph{Amazon VPC pricing}.}
\end{frame}

\begin{frame}{Topologie réseau type}
  \centering
  \includegraphics[height=.8\textheight]{vpc-topologie}
\end{frame}

\begin{frame}{Notation CIDR}
  \renewcommand{\arraystretch}{1.4}
  \begin{tabular}{@{}l l r@{}}
    \toprule
    \textbf{Notation} & \textbf{Signification} & \textbf{Adresses} \\ \midrule
    \texttt{10.0.0.0/16} & bloc d'un VPC & 65 536 \\
    \texttt{10.0.1.0/24} & bloc d'un sous-réseau & 256 (251 utilisables sur AWS) \\
    \texttt{203.0.113.25/32} & une seule adresse & 1 \\
    \texttt{0.0.0.0/0} & toutes les adresses IPv4 & \\
    \texttt{::/0} & toutes les adresses IPv6 & \\ \bottomrule
  \end{tabular}
  \bigskip
  \begin{itemize}
    \item AWS réserve 5 adresses par sous-réseau
  \end{itemize}
\end{frame}

\begin{frame}{Security Groups}
  \begin{itemize}
    \setlength{\itemsep}{.45em}
    \item Pare-feu virtuel attaché à l'interface réseau d'une ressource (instance, base RDS...)
    \item Règles : protocole, port, source (CIDR ou autre Security Group)
    \item Uniquement des autorisations
    \item Par défaut : entrant interdit, sortant autorisé
    \item \textbf{Avec état} (\emph{stateful}) : la réponse à une connexion autorisée est autorisée
    \item Modifications appliquées immédiatement
    \item Plusieurs Security Groups par ressource : règles cumulées
  \end{itemize}
  \lien{https://docs.aws.amazon.com/vpc/latest/userguide/vpc-security-groups.html}
\end{frame}

\begin{frame}{Security Group d'un serveur web}
  \centering
  \includegraphics[width=.95\textwidth]{security-group}
\end{frame}

\begin{frame}{Security Groups et NACL}
  \renewcommand{\arraystretch}{1.4}
  \begin{tabularx}{\textwidth}{@{}>{\raggedright\arraybackslash}p{3.4cm}>{\raggedright\arraybackslash}X>{\raggedright\arraybackslash}X@{}}
    \toprule
    & \textbf{Security Group} & \textbf{Network ACL} \\ \midrule
    Portée & interface réseau (instance) & sous-réseau \\
    Règles & autorisations seulement & autorisations et interdictions \\
    État & avec état & sans état : trafic retour à autoriser \\
    Évaluation & toutes les règles & par numéro de règle, la première qui correspond \\
    Par défaut & entrant refusé, sortant autorisé & NACL par défaut : tout autorisé \\ \bottomrule
  \end{tabularx}
\end{frame}

\begin{frame}{Bonnes pratiques réseau}
  \begin{itemize}
    \setlength{\itemsep}{.5em}
    \item N'ouvrir que les ports nécessaires
    \item SSH (22) limité à une adresse précise, jamais \texttt{0.0.0.0/0}
    \item Bases de données dans des sous-réseaux privés
    \item Source = Security Group plutôt que plage d'adresses entre niveaux d'une application
    \item Session Manager (Systems Manager) : administration sans port SSH ouvert
  \end{itemize}
  \bigskip
  \attention L'espace IPv4 complet peut être balayé en moins de 45 minutes depuis une seule machine (ZMap)
  \source{Durumeric et al., « ZMap: Fast Internet-Wide Scanning and its Security Applications », USENIX Security 2013.}
\end{frame}

\diapotp{2}{Utilisateurs, droits et pare-feu}{%
  \item Stratégie sur mesure, groupe et utilisateur au moindre privilège
  \item Tests dans la console et dans CloudShell
  \item Simulateur de stratégies, refus explicite
  \item Security Group du serveur web}

